\section{Agent 系统设计蓝图：从研究概念到可运行产品}

前面几章分别讨论语言、任务环境、reward 和 multi-agent，本章把它们合成一个可设计的系统蓝图。一个 Agent 产品通常不是单个 prompt，而是一套循环：用户目标进入系统，系统把目标转成任务状态，模型生成计划，工具执行动作，环境返回观察，验证器判断结果，记忆层保存经验，最后再进入下一轮。

\begin{importantbox}{Agent 系统最小闭环}
\[
\text{Goal} \rightarrow \text{State} \rightarrow \text{Plan} \rightarrow \text{Action} \rightarrow \text{Observation} \rightarrow \text{Reward/Verification} \rightarrow \text{Memory}
\]
其中，Goal 是用户目标，State 是当前任务状态，Plan 是模型生成的计划，Action 是工具或代码执行，Observation 是环境反馈，Reward/Verification 是评价信号，Memory 是长期经验和上下文。
\end{importantbox}

\begin{knowledgebox}{设计蓝图：每一层问什么}
\begin{tabularx}{\textwidth}{p{0.20\textwidth}p{0.34\textwidth}X}
\toprule
层级 & 设计问题 & 常见失败 \\
\midrule
目标层 & 用户到底要什么？成功条件是什么？ & 目标含糊，Agent 做了看似有用但无关的事。 \\
状态层 & 当前有哪些文件、工具、权限和约束？ & 上下文缺失，动作建立在错误状态上。 \\
计划层 & 是否需要拆任务、并行或询问用户？ & 计划太长、太脆弱，遇到异常不会改。 \\
行动层 & 哪些动作能安全执行，哪些必须确认？ & 越权、误删、调用错误工具。 \\
验证层 & 如何知道任务真的完成了？ & 只看表面输出，不验证真实结果。 \\
记忆层 & 哪些经验应该保留到下次？ & 记忆污染、隐私风险、过期信息。 \\
\bottomrule
\end{tabularx}
\end{knowledgebox}

\begin{warningbox}{从 demo 到产品的距离}
Agent demo 可以靠一次顺利轨迹打动人；Agent 产品必须处理失败恢复、权限、日志、验证、用户中断、长期记忆和成本。下半场的系统能力，很大一部分就在这些“无聊但必要”的层里。
\end{warningbox}

\subsection{本章小结}

Agent 系统设计的核心不是把模型包一层 UI，而是把目标、状态、计划、行动、观察、验证和记忆闭合起来。只有闭环稳定，Agent 才能从惊艳演示变成可靠产品。

